%قبل از شروع مسئله ترجیحاً چیزی ننویسید. اما اگر خواستید بنویسید، در قالب لیستی (فقرات و شمارش) بنویسید.
%برای اضافه کردن نمره، بعدی از شروع‌مسئله، [x] را اضافه کنید و به جای x  نمره مورد نظر را قرار دهید. برای زیرمسئله‌ها هم به همین صورت است.
%\شروع‌مسئله[10]
\شروع‌مسئله

% صورت اصلی مسئله را در این‌جا بنویسید.

یک آرایه بیتی (صفر و یکی) $m$ تایی فضای بسیار کمتری از آرایه‌ای از $m$ پوینتر می‌گیرد. در صورتی که عناصری که داریم، داده اضافی همراه خود نداشته باشند و صرفا مقادیری بین $0$ تا $m-1$ داشته باشند، نشان دهید چگونه روش آدرس‌دهی مستقیم را می‌توان با آرایه بیتی پیاده سازی کرد.
\solution{
این کار به شکل بسیار ساده قابل انجام است. فرض کنید یک آرایه بیتی $20$ تایی داریم و مقادیر $2,4,6,10,15$ در آن قرار گرفته باشند. این را به راحتی با مقدار زیر می‌توان نمایش داد:
$$00101010001000010000$$

بر این اساس، اگر بدانیم فضای حالت ما به اندازه کافی کوچک است، روش آدرس‌دهی مستقیم از طریق این توابع قابل پیاده سازی خواهد بود؛ توجه کنید که $V$ همان آرایه بیتی ما است:


\begin{latin}
	\begin{algorithm}[H]
		\begin{algorithmic}
			
			\Procedure{$Bitmap\_Search(V,k)$}{}
			\If {$V[k] \neq 0$}
			\Return k
			\EndIf
			\Return Null
			\EndProcedure
			
			
			\Procedure{$Bitmap\_Insert(V,k)$}{}
			\State $V[k] = 1$
			\EndProcedure
			
		\Procedure{$Bitmap\_Delete(V,k)$}{}
		\State $V[k] = 0$
		\EndProcedure
			
		\end{algorithmic}
	\end{algorithm}	
	
\end{latin}




}



%%%%%%%%%%%%%%%%%%%%%%%%%%%%%%%%%%%%%%%%%%%%%%%%%%%%%
% برخی دستورات مفید برای تمرین‌ها. %
%%%%%%%%%%%%%%%%%%%%%%%%%%%%%%%%%%%%%%%%%%%%%%%%%%%%%

% فرمول‌های نمونه:
%$f(x) = \Theta(x)$
%$T(x, y) \leq \log(x) + 2^y + \binom{2n}{n}$
% {\tt code\_function}


% لیست موارد بدون شماره‌گذاری:
%\شروع{فقرات}
%    \فقره مورد اول 
%        \شروع{فقرات}
%            \فقره مورد اول داخلی 
%                \شروع{فقرات}
%                    \فقره اولین مورد خیلی داخلی 
%                    \فقره دومین مورد خیلی داخلی 
%                \پایان{فقرات}
%            \فقره مورد دوم داخلی 
%        \پایان{فقرات}
%    \فقره مورد دوم 
%\پایان{فقرات}

% لیست موارد شماره‌گذاری شده:
%\شروع{شمارش}
%    \فقره یک مورد 
%    \فقره یک مورد دیگر 
%    \فقره باز هم یک مورد دیگر
%\پایان{شمارش}

% برای نوشتن یک فرمول چند خطی که یک شماره بخوره. برای حالت‌های دیگر به سایت زیر مراجعه کنید:
% https://www.overleaf.com/learn/latex/Aligning_equations_with_amsmath
%\begin{align} 
%    \begin{split}
%        (x+y)^3 &= (x+y)^2(x+y) \\
%                &= (x^2+2xy+y^2)(x+y) \\
%                &= (x^3+2x^2y+xy^2) + (x^2y+2xy^2+y^3) \\
%                &= x^3+3x^2y+3xy^2+y^3
%    \end{split}                                 
%\end{align}

% نوشتن ماتریس و آرایه دوبعدی
%\begin{align}
%    A = 
%    \begin{bmatrix}
%        A_{11} & A_{21} \\
%        A_{21} & A_{22}
%    \end{bmatrix}
%\end{align}

% اضافه کردن تصویر با شرح. برای اضافه کردن تصویر بدون شرح، فقط دستور درج‌تصویر را استفاده کنید.
%\شروع{شکل}[h!]
%\تنظیم‌ازوسط
%  \درج‌تصویر[پهنا=0.5\پهنای‌متن]{img.jpg}
%\شرح{شرح تصویر}
%\برچسب{تصویر:۱}
%\پایان{شکل}
%\پایان{زیرمسائل}

%کد زیر را ببینید.
%%\lstinputlisting{source_filename.py}
%\begin{latin}
%\begin{lstlisting}
%# Import the modules
%import sys
%import random
%
%ans = True
%
%while ans:
%question = raw_input("Ask the magic 8 ball a question: (press enter to quit) ")
%
%answers = random.randint(1,8)
%
%if question == "":
%sys.exit()
%
%elif answers == 1:
%print "It is certain"
%
%elif answers == 2:
%print "Outlook good"
%
%elif answers == 3:
%print "You may rely on it"
%
%elif answers == 4:
%print "Ask again later"
%
%elif answers == 5:
%print "Concentrate and ask again"
%
%elif answers == 6:
%print "Reply hazy, try again"
%
%elif answers == 7:
%print "My reply is no"
%
%elif answers == 8:
%print "My sources say no"
%\end{lstlisting}
%\end{latin}
